\documentclass[11pt, letterpaper]{article}
\input{../../homework.tex}
\usepackage{titlepageBU}
\title{Homework 2}
\courseID{PY 355}
\courseSection{A1}
\begin{document}
\makeproblem
\section*{Problem 1}
\begin{enumerate}
	\item Exponent rule
		\[
			\int_{0}^{1} a^y \,\mathrm{d} y= \left.\frac{a^y}{\ln a}  \right|_0^1 = \frac{a-1}{\ln a} 
		.\]
	\item U-substitution, $u=ay$, $du=adx$, $v=-ay$, $dv=-vdx$.
		\[
			\int_{0}^{x} \cosh y \,\mathrm{d} y=\int_{0}^{x} \frac{e^{ax} +e^{-ax} }{2}  \,\mathrm{d} y= \frac{1}{2} \left( \int_{0}^{x} e^{ay}  \,\mathrm{d} y+\int_{0}^{x} e^{-ay}  \,\mathrm{d} y \right)
		\]
		\[
			=\frac{1}{2a} \left( \int_{0}^{ax} e^u \,\mathrm{d} u-\int_{0}^{-ax} e^v \,\mathrm{d} v \right) =\frac{1}{2a} \left.e^u\right|_{0}^{ax}-\left. e^{v}\right|_0^{-ax} = \frac{1}{2a} \left( e^{ax} -1-e^{-ax} +1 \right) 
		\]
		\[
			=\frac{1}{a} \frac{e^{ax} -e^{-ax} }{2} =\frac{1}{a} \sinh(ax)
		.\]
	\item Derivative of $\ln x$, u substitution $u=ay+b$, $du=1/a\,dx$.
		\[
			\int_{0}^{x} \frac{1}{ay+b}  \,\mathrm{d} y= \int_{0}^{x} \left( ax+b \right) ^{-1}  \,\mathrm{d} y=\int_{0}^{x} \frac{1}{a} u^{-1}  \,\mathrm{d} u= \frac{1}{a} \ln(ay+b) \bigg|_0^x = \frac{1}{a} \left( \ln (ax+b)-\ln b \right) 
		.\]
	\item IBP, let $g=y^2$, $g'=2y$, $f'=\cos (2y)$, $f=1/2 \cos (2y)$.
		\[
			\int_{0}^{x} y^2 \cos (2y) \,\mathrm{d} y=\frac{y^2}{2} \sin (2y)\bigg|_0^x -\int_{0}^{x} \frac{2y}{2} \sin (2y) \,\mathrm{d} y =\frac{x^2}{2} \sin (2x) -\int_{0}^{x} y \sin (2y) \,\mathrm{d} y
		.\]
		IBP, let $g=y$, $g'=1$, $f'=\sin (2y)$, $f=-1/2 \cos (2y)$.
		\[
			=\frac{x^2}{2} \sin (2x)+\frac{y}{2} \cos (2y)\bigg|_0^x-\int_{0}^{x} \frac{1}{2} \cos (2y) \,\mathrm{d} y=\frac{x^2}{2} \sin (2x)+\frac{x}{2} \cos (2x)-\frac{1}{4} \sin (2y)\bigg|_0^x
		\]
		\[
			=\frac{1}{2} \left( x^2 \sin (2x)+ x \cos (2x)- \frac{1}{2} \sin (2x) \right) 
		.\]
	\item U-substitution, let $u=-x^4$, $du=-4x^3 dx$.
		\[
			\int_{0}^{z} x^3e^{-x^4}  \,\mathrm{d} x=\int_{0}^{-z^4} -\frac{1}{4} e^u \,\mathrm{d} u=-\frac{1}{4} e^u \bigg|_0^{-z^4}=\frac{1}{4} \left(-e^{-z^4} +1\right)
		.\]
	\item Starting with a proof that $\cosh^2 x-\sinh^2 x=1$.
		\[
			\left( \frac{e^{x} +e^{-x} }{2}  \right)^2-\left( \frac{e^{x} -e^{-x} }{2}  \right) ^2=\left( \frac{e^{2x} +2+e^{-2x} }{4}  \right) -\left( \frac{e^{2x} -2+e^{-2x} }{4}  \right) =\frac{4}{4} =1
		.\]
		U-substitution, $u=\sinh x$, $du=\cosh x$.
		\[
			\int_{0}^{z} \cosh^3 x \,\mathrm{d} x=\int_{0}^{z} \left( 1+\sinh^2x \right) \cosh x \,\mathrm{d} x=\int_{0}^{\sinh z} 1+u^2 \,\mathrm{d} u=\left( u+\frac{1}{3} u^3 \right)\bigg|_0^{\sinh z}
		\]
		\[
			=\sinh z+ \frac{1}{3} \sinh^3 z
		.\]
	\item Trig substitution, let $x=a\sin \theta $. Then $dx=a \cos \theta d \theta $. Also if $x=z$ then $\theta = \arcsin \frac{z}{a} $, and if $x=0$ then $\theta =0$.
		\[
			\int_{0}^{z} \frac{1}{\sqrt{a^2-x^2} }  \,\mathrm{d} x=\int_{0}^{\arcsin z/a} \frac{a \cos \theta }{\sqrt{a^2\left( 1-\sin^2 \theta  \right) } }  \,\mathrm{d} \theta =\int_{0}^{\arcsin z/a} \frac{a \cos \theta }{a\sqrt{\cos ^2\theta } }  \,\mathrm{d} \theta 
		\]
		\[
			=\int_{0}^{\arcsin z/a} \frac{a \cos \theta }{a \cos \theta }  \,\mathrm{d} \theta =\int_{0}^{\arcsin z/a}  \,\mathrm{d} \theta =\arcsin \frac{z}{a} 
		.\]
		This works because we take $\left| z \right| <a$, so $\left| z/a \right| <1$, and $\arcsin $ is well-defined on this interval.
\end{enumerate}
\section*{Problem 2}
Let $g=\ln x$, $g'=1/x$, $f'=1$, $f=x$. Then
\[
	\int \ln x \,\mathrm{d} x=x \ln x-\int \frac{x}{x}  \,\mathrm{d} x=x \ln x-\int  \,\mathrm{d} x=x \ln x-x+C
.\]
\section*{Problem 3}
Isn't this basically just $\Gamma (n+1)$? First, we show $F(n)=nF(n-1)$ for $n\geq 1$. Let $g=x^n$ and $f'=e^{-x} $. Then $g'=nx^{n-1} $ and $f=-e^{-x} $. We have
\[
	F(n)=\int_{0}^{\infty } x^ne^{-x}  \,\mathrm{d} x=-x^ne^{-x}\bigg|_0^{\infty }+\int_{0}^{\infty } nx^{n-1} e^{-x}  \,\mathrm{d} x=\left( \lim_{x \to \infty} \frac{-x^n }{e^x}  \right) +nF(n-1)
.\]
It suffices to show the limit is 0. Note that if $n-1\geq 1$, we have
\[
	\lim_{x \to \infty} \frac{-x^n }{e^{x} } =-\frac{\infty }{\infty } 
.\]
Therefore, we can apply L'Hopital's rule until the exponent is $0$. We then have
\[
	\lim_{x \to \infty} \frac{-x^n }{e^x} =\lim_{x \to \infty} \frac{-n!}{e^x} =0
\]
because $n$ is finite, so $n!$ is finite.

Next, we will compute $F(0)$.
\[
	F(0)=\int_{0}^{\infty } e^{-x}  \,\mathrm{d} x=-e^{-x} \bigg|_0^{\infty }=0-(-1)=1
.\]
Therefore, $F(0)=1$.

Now we will show $F(n)=n!$ for $n\in\mathbb{N} $ using the method of induction. Since $F(n)=nF(n-1)$, and $F(0)=1$, we have
\[
	F(1)=1\cdot F(0)=1
.\]
Therefore, $F(1)=1=1!$. Now suppose that $F(n)=n!$ for some arbitrary $n\in\mathbb{N} $. It follows that
\[
	F(n+1)=(n+1)F(n)=(n+1)n! =(n+1)!
.\]
Therefore, $F(n+1)=(n+1)!$, and the statement is proven.

The integral defines $0! = 1$ because $F(0)=1$.
\section*{Problem 4}
\begin{enumerate}
	\item Let $u=\sqrt{a} x$. Then
\[
	I_0(a)=\int_{-\infty }^{\infty } e^{-ax^2}  \,\mathrm{d} x=\frac{1}{\sqrt{a} } \int_{-\infty }^{\infty } e^{-u^2}  \,\mathrm{d} u=\frac{1}{\sqrt{a} } I_0(1)
.\]
	\item Since $n$ is odd, let $n=2m+1$ for some $m\in\mathbb{N}$. We have
		\[
			I_n(a)=\int_{-\infty }^{\infty } x^{2m+1} e^{-ax^2}  \,\mathrm{d} x=\int_{-\infty }^{\infty } x\left( \frac{2a}{2a} x^2 \right)^{m}  e^{-ax^2}  \,\mathrm{d} x
		.\]
		Let $u=ax^2$. Then $du=2axdx$. Note that applying this $u$-substitution will cause both bounds to now be $\infty $ instead of $\pm \infty $, so I will break the integral into two parts to make stuff work. We have
		\[
			=\frac{1}{\left( 2a \right)^{m+1} } \left( \int_{0}^{\infty } u^me^{-u}  \,\mathrm{d} u+\int_{-\infty }^{0} u^me^{-u}  \,\mathrm{d} u \right)
		.\]
		\begin{align*}
			\int_{-\infty }^{\infty } x\left( \frac{2a}{2a} x^2 \right)^{m}  e^{-ax^2}  \,\mathrm{d} x&=\int_{-\infty }^{0} x\left( \frac{2a}{2a} x^2 \right)^m e^{-ax^2}  \,\mathrm{d} x+\int_{-\infty }^{0} x\left( \frac{2a}{2a} x^2 \right)^m e^{-ax^2}  \,\mathrm{d} x\\
														  &=\int_{-\infty }^{0} \frac{1}{\left( 2a \right)^{m+1} } u^{m} e^{-u}  \,\mathrm{d} u+\int_{0}^{\infty } \frac{1}{\left( 2a \right)^{m+1} } u^{m} e^{-u}  \,\mathrm{d} u\\
														  &=\frac{1}{\left( 2a \right) ^{m+1} } \left( -\int_{0}^{\infty } u^m e^{-u}  \,\mathrm{d} u+\int_{0}^{\infty } u^m e^{-u}  \,\mathrm{d} u \right)
		.\end{align*}
		Since both integrals are now the same, they cancel out, and $I_n(a)=0$. Interestingly, note both of these integrals are $F(m)$ as defined in problem 3.
	\item We use proof by the method of induction. Let $p=1$. Then
		\[
			(-1) \frac{\mathrm{d}I_0(a)}{\mathrm{d}a} = -\int_{-\infty }^{\infty } \frac{\mathrm{d}}{\mathrm{d}a} e^{-ax^2}  \,\mathrm{d} x=-\int_{-\infty }^{\infty } \left( -x^2 \right) e^{-ax^2}  \,\mathrm{d} x=\int_{-\infty }^{\infty } x^2e^{-ax^2}  \,\mathrm{d} =I_2(a)
		.\]
		Now suppose the statement holds for some arbitrary $p\in\mathbb{N} $. We have
		\[
			(-1)^{p+1} \frac{\mathrm{d}^{p+1} I_0(a)}{\mathrm{d}a^{p+1} } =-\int_{-\infty }^{\infty } \frac{\mathrm{d}}{\mathrm{d}a} x^{2p} e^{-ax^2}  \,\mathrm{d} x=-\int_{-\infty }^{\infty } x^{2p} \left( -x^2 \right) e^{-ax^2}  \,\mathrm{d} x
		\]
		\[
			=\int_{-\infty }^{\infty } x^{2(p+1)} e^{-ax^2}  \,\mathrm{d} x=I_{2(p+1)} (a)
		.\]
		Therefore,
		\[
			(-1)^{p} \frac{\mathrm{d}^pI_0(a)}{\mathrm{d}a^p} =I_{2p} (a)
		.\]
	\item The problem does not ask for proof, only for the pattern.
		\[
			I_0(a)=\sqrt{\frac{\pi }{a} } 
		.\]
		We know this from part 1 and commentary in the problem set.
		\[
			I_2(a)=-\frac{\mathrm{d}I_0(a)}{\mathrm{d}a} =\frac{\mathrm{d}}{\mathrm{d}a} \sqrt{\pi } a^{-1/2} =\frac{1}{2} \sqrt{\frac{\pi }{a^3} } 
		.\]
		\[
			I_4(a)=-\frac{\mathrm{d}I_2(a)}{\mathrm{d}a} =-\frac{\mathrm{d}}{\mathrm{d}a}\frac{1}{2}  \sqrt{\pi } a^{-3/2} =\frac{3}{2} \cdot \frac{1}{2} \sqrt{\frac{\pi }{a^{5} } } 
		.\]
		In general, the formula is
		\[
			I_{2n} (a)=\frac{1}{2^n} \left( \prod_{i=0}^{n-1} 2i+1 \right) \sqrt{\frac{\pi }{a^{2n+1} } }=\frac{(2n)!}{2^{2n} n!} \sqrt{\frac{\pi }{a^{2n+1} } } 
		.\]
		This second representation uses \href{https://math.stackexchange.com/a/69163}{this identity}.
\end{enumerate}
\section*{Problem 5}
\begin{enumerate}
	\item The argument is exactly the same, so I will show less work. We will prove by induction. Let $p=1$. then
		\[
			-\frac{\mathrm{d}\widetilde{I}_1(a)}{\mathrm{d}a} =-\int_{0}^{\infty } \frac{\mathrm{d}}{\mathrm{d}a} xe^{-ax^2}  \,\mathrm{d} x=\int_{0}^{\infty } x^3 e^{-ax^2}  \,\mathrm{d} x=\widetilde{I} _3(a)=\widetilde{I} _{2+1} (a)
		.\]
		Now suppose the statement holds for $p\in\mathbb{N} $ arbitary, and consider $p+1$. We have
		\[
			- \frac{\mathrm{d}^{p+1} \widetilde{I}_1(a)}{\mathrm{d}a^{p+1} } =-\int_{0}^{\infty } \frac{\mathrm{d}}{\mathrm{d}a} x^{2p+1} e^{-ax^2}  \,\mathrm{d} x=\int_{0}^{\infty } x^{2p+3} e^{-ax^2}  \,\mathrm{d} x=\int_{0}^{\infty } x^{2(p+1)+1} e^{-ax^2}  \,\mathrm{d} x
		\]
		\[
			=\widetilde{I}_{2(p+1)+1} (a)
		.\]
	By the method of induction, the theorem is proven.

	\item Let $u=ax^2$. Then $du=2axdx$. Recall the definition of $F(n)$ from problem 3, and that $F(0)=1$. We have
		\[
			\widetilde{I}_1(a)=\int_{0}^{\infty } xe^{-ax^2}  \,\mathrm{d} x= \frac{1}{2a}\int_{0}^{\infty } e^{-u}  \,\mathrm{d} u=\frac{1}{2a} F(0)=\frac{1}{2a} 
		.\]
\end{enumerate}

\end{document}
